\documentclass[11pt]{article}
% cs250preamble.tex --- shared preamble for CS 250 course notes
%
% This file is \input by main.tex and by each individual module
% file so that each module can still compile standalone. Any
% package or custom-environment change should be made here, not
% in the individual module files.

\usepackage[T1]{fontenc}
\usepackage[margin=1in]{geometry}
\usepackage{amsmath, amssymb, amsthm}
\usepackage{enumitem}
\usepackage{hyperref}
\usepackage{booktabs}
\usepackage{listings}
\usepackage{xcolor}
\usepackage{tikz}
\usetikzlibrary{shapes.geometric, shapes.gates.logic.US, shapes.gates.logic.IEC, arrows.meta, positioning, calc, trees}
\usepackage{bussproofs}
\usepackage{mdframed}

\lstset{
  basicstyle=\ttfamily\small,
  frame=single,
  breaklines=true,
  columns=fullflexible,
  mathescape=true
}

\theoremstyle{definition}
\newtheorem{theorem}{Theorem}[section]
\newtheorem{definition}[theorem]{Definition}
\newtheorem{example}[theorem]{Example}

\hypersetup{
  colorlinks=true,
  linkcolor=blue!60!black,
  urlcolor=blue!60!black
}

% Key-takeaway box: used at the end of major sections and at the end of
% each module to summarise the main points. Expects \item bullets inside.
\newmdenv[
  linewidth=0.8pt,
  linecolor=blue!40!black,
  backgroundcolor=blue!3,
  skipabove=8pt,
  skipbelow=8pt,
  innertopmargin=7pt,
  innerbottommargin=7pt,
  innerleftmargin=9pt,
  innerrightmargin=9pt
]{keytakeawaybox}
\newenvironment{keytakeaway}{%
  \begin{keytakeawaybox}%
  \noindent\textbf{Key takeaways.}%
  \begin{itemize}[leftmargin=1.3em, topsep=2pt, itemsep=1pt, label=\textbullet]%
}{%
  \end{itemize}%
  \end{keytakeawaybox}%
}

% Summary/looking-ahead environment: used once at the end of each module
% to wrap up what the module built and point forward.
\newmdenv[
  linewidth=0.8pt,
  linecolor=gray!60,
  backgroundcolor=gray!5,
  skipabove=8pt,
  skipbelow=8pt,
  innertopmargin=7pt,
  innerbottommargin=7pt,
  innerleftmargin=9pt,
  innerrightmargin=9pt
]{summarybox}

% Sidebar environment: used for tangential asides---historical notes,
% pointers to other areas of math/CS, cross-paradigm remarks---that
% enrich the main narrative without being essential to it. Takes one
% argument: the sidebar's title.
\newmdenv[
  linewidth=0.6pt,
  linecolor=gray!60,
  backgroundcolor=gray!4,
  skipabove=8pt,
  skipbelow=8pt,
  innertopmargin=6pt,
  innerbottommargin=6pt,
  innerleftmargin=9pt,
  innerrightmargin=9pt
]{sidebarbox}
\newenvironment{sidebar}[1]{%
  \begin{sidebarbox}%
  \noindent\textbf{Sidebar: #1.}\par\medskip%
}{%
  \end{sidebarbox}%
}

% Reading-map environment: used at the top of each numbered module to
% indicate Epp coverage, prerequisites, relevant intermezzi, and
% suggested practice problems. Expects four \rmentry{label}{body}
% items inside (or arbitrary paragraphs).
\newmdenv[
  linewidth=0.8pt,
  linecolor=black!55,
  backgroundcolor=yellow!4,
  skipabove=10pt,
  skipbelow=14pt,
  innertopmargin=9pt,
  innerbottommargin=9pt,
  innerleftmargin=11pt,
  innerrightmargin=11pt
]{readingmapbox}
\newcommand{\rmentry}[2]{\par\noindent\textbf{#1}\hspace{0.4em}#2\par}
\newenvironment{readingmap}{%
  \begin{readingmapbox}%
  \noindent{\bfseries\large Reading map}%
  \vspace{2pt}%
  \setlength{\parskip}{4pt plus 1pt}%
}{%
  \end{readingmapbox}%
}


\title{CS 250 --- Cumulative Review 2 (after Module 6)}
\author{}
\date{}

\begin{document}
\maketitle

\section*{The point of this review}

We're now halfway through the first term, and the toolkit has grown considerably. In Module~4 we added \emph{predicate logic} and a first taste of induction on $\mathbb{N}$; in Module~5 we learned to write \emph{informal proofs} (direct, counterexample, existence, and the group axioms); in Module~6 we introduced \emph{modular arithmetic} and \emph{proof by cases} on residue classes. Adding those to the set-and-propositional-logic foundation from Modules 1--3 gives us enough machinery to prove real, non-trivial theorems---the kind that appear in number theory textbooks.

In this review chapter we pick one classical result, prove it carefully, and then trace through the proof to highlight how every module's tools contribute. The result is:

\begin{theorem}
For every integer $n$, the quantity $n^3 - n$ is divisible by $6$.
\end{theorem}

This is the right level of difficulty for a mid-term review: it is clearly non-trivial (you have to combine at least two separate arguments to get the factor of $6$), but it is also clean enough that the proof fits on one page. Let's go.

\section*{The proof}

We first rewrite the quantity in a more suggestive form:
\[
  n^3 - n = n(n^2 - 1) = n(n-1)(n+1) = (n-1) \cdot n \cdot (n+1).
\]
So $n^3 - n$ is the product of three consecutive integers.

To show $6 \mid n^3 - n$, it suffices (by a consequence of the Chinese Remainder Theorem, since $\gcd(2, 3) = 1$) to show that both $2$ and $3$ divide $n^3 - n$. We'll handle each divisibility separately.

\subsection*{Step 1: $2 \mid n^3 - n$}

Among any two consecutive integers, one is even. So among $n-1$, $n$, $n+1$, at least one is even---in fact, \emph{at least} one (and possibly two) are even. Therefore the product $(n-1) \cdot n \cdot (n+1)$ is even, i.e., $2 \mid n^3 - n$.

If you want a more formal version of this: proceed by cases on the parity of $n$ (the two-residue-class partition of $\mathbb{Z}$ we get from the quotient-remainder theorem mod $2$):

\emph{Case 1: $n$ is even.} Then $n = 2k$ for some integer $k$, and $(n-1) \cdot n \cdot (n+1) = 2k \cdot (n-1)(n+1) = 2 \cdot [k(n-1)(n+1)]$. The bracketed quantity is an integer, so $2$ divides the product.

\emph{Case 2: $n$ is odd.} Then $n - 1$ is even, so $n - 1 = 2j$ for some integer $j$, and the product becomes $2j \cdot n \cdot (n+1) = 2 \cdot [j \cdot n \cdot (n+1)]$, which is even.

Either way, the product is even.

\subsection*{Step 2: $3 \mid n^3 - n$}

This is the more interesting half. We use \emph{proof by cases on $n \bmod 3$}: by the quotient-remainder theorem, every integer $n$ satisfies exactly one of $n \equiv 0, 1, 2 \pmod 3$.

\emph{Case 1: $n \equiv 0 \pmod 3$.} Then $3 \mid n$, so $3 \mid n(n-1)(n+1)$.

\emph{Case 2: $n \equiv 1 \pmod 3$.} Then $n - 1 \equiv 0 \pmod 3$, so $3 \mid (n-1)$, so $3 \mid n(n-1)(n+1)$.

\emph{Case 3: $n \equiv 2 \pmod 3$.} Then $n + 1 \equiv 0 \pmod 3$, so $3 \mid (n+1)$, so $3 \mid n(n-1)(n+1)$.

In every case, $3$ divides the product.

\subsection*{Combining: $6 \mid n^3 - n$}

We've shown that $2 \mid n^3 - n$ and $3 \mid n^3 - n$. Since $\gcd(2, 3) = 1$, the two divisibilities combine: if integers $a$ and $b$ are coprime and both divide $m$, then $ab$ divides $m$.

This is the mini-CRT fact we noted in Module~6. Here's a short direct argument: by B\'ezout's identity (coming in Module~7, but elementary enough to give here), since $\gcd(2, 3) = 1$ there exist integers $x, y$ with $2x + 3y = 1$. Multiplying by $m$, we get $m = 2xm + 3ym$. If $2 \mid m$ (so $m = 2m'$) and $3 \mid m$ (so $m = 3m''$), then $m = 2xm + 3ym = 2x(3m'') + 3y(2m') = 6(xm'' + ym')$, which is divisible by $6$.

So $6 \mid n^3 - n$ for every integer $n$. \hfill$\square$

\section*{Tracing the toolkit}

Now let's go back and count which modules contributed what to that proof.

\begin{itemize}
  \item \textbf{Module 1} (sets, functions, proof-writing) gave us the disciplined prose style: state the goal, introduce variables, justify each step, conclude clearly. Every ``Let $n = 2k$ for some integer $k$'' and every ``By definition of divisibility\ldots'' is a habit from Module~1.
  \item \textbf{Module 2} (propositional logic) sat quietly in the background. Every case analysis we did was justified by a tautology of propositional logic---specifically, the disjunction $(n \equiv 0 \pmod 3) \vee (n \equiv 1 \pmod 3) \vee (n \equiv 2 \pmod 3)$, combined with $\vee$-elimination.
  \item \textbf{Module 3} (proof systems) is what makes our casual use of $\vee$-elimination rigorous. The ``proof by cases'' we did in Step~2 is formally an application of the $\vee$-elimination rule from Module~3; each case branch is a subproof. The Gentzen intermezzo sketches how to draw the full derivation as a tree.
  \item \textbf{Module 4} (predicate logic and induction) gave us the quantifier vocabulary. ``Every integer $n$'' is $\forall n : \mathbb{Z}$; ``some integer $k$ with $n = 2k$'' is $\exists k : \mathbb{Z}$. When we write ``let $n$ be an arbitrary integer,'' we are invoking $\forall$-introduction. When we write ``$n = 2k$ for some integer $k$,'' we are using $\exists$-elimination to extract the witness from the hypothesis ``$n$ is even.''
  \item \textbf{Module 5} (direct proof, groups) provided the overall proof style: unpack the definitions (``even means $n = 2k$''), do the algebra, repack the result (``hence $2 \mid n^3 - n$''). The uniqueness arguments in Module~5 trained the habit of asking ``is the thing I just chose arbitrary?''
  \item \textbf{Module 6} (modular arithmetic and proof by cases) is the star of this proof. The case split on $n \bmod 3$ uses the quotient-remainder theorem to guarantee the three cases are exhaustive, and the final CRT-flavored combination step uses precisely the mini-CRT we highlighted in Module~6. The entire shape of Step~2 is Module-6 shape.
\end{itemize}

Every module contributed, and dropping any one of them would make the proof harder (or would force us to use a different approach). The moral is that the modules aren't independent silos: each one adds a layer of tooling that sits on top of the previous ones, and a real proof freely draws on all of them at once.

\section*{A variant to chew on}

Here's a closely related claim you can try yourself: \emph{for every integer $n$, the quantity $n^5 - n$ is divisible by $30$.} (Note: $30 = 2 \cdot 3 \cdot 5$.) You'll need cases on $n \bmod 5$ and a similar coprime-combination step. Try it. If you can do it, you've earned the first half of the term.

A slightly harder variant: \emph{for every prime $p$ and every integer $n$, $p \mid n^p - n$.} This is Fermat's Little Theorem, and it generalizes both of the above. The standard proof uses induction (Module~8) plus the binomial theorem, so you'll want to revisit this once you've done Module~8.

\section*{Where we're going}

Modules 7--10 take the toolkit two more steps forward:
\begin{itemize}
  \item Module~7 adds \emph{indirect proof} (contradiction and contrapositive), the \emph{Euclidean algorithm} (which makes B\'ezout's identity computational), and a first taste of \emph{Hoare logic} for program correctness.
  \item Modules 8 and 9 revisit \emph{induction} in a bigger way: summation identities, strong induction, recurrence solving.
  \item Module~10 generalizes induction to arbitrary recursive data types (\emph{structural induction}), completing the arc that started with BNF definitions in Module~2.
\end{itemize}

After Module~10, the third cumulative review chapter ties everything together.

\end{document}
